\chapter{Funkcje mierzalne}
\label{chap:funkcje-mierzalne}

\section{Równoważne warunki mierzalności}
\label{sec:warunki-mierzalnosci-funkcji}

W tym rozdziale $(X,\mathcal{F})$ jest przestrzenią mierzalną.
Dla funkcji $f\colon X\to\R$ i zbioru $B\subseteq\R$ oznaczamy
\[
  f^{-1}(B)=\{x\in X:f(x)\in B\}.
\]
Jest to przeciwobraz zbioru; zapis nie wymaga odwracalności funkcji.

\begin{twierdzenie}[Równoważne warunki]
\label{tw:warunki-mierzalnosci-funkcji}
Dla funkcji $f\colon X\to\R$ następujące warunki są równoważne:
\begin{enumerate}
  \item dla każdego $a\in\R$ zbiór
    $\{x\in X:f(x)>a\}=f^{-1}((a,+\infty))$ należy do $\mathcal{F}$;
  \item dla każdego $a\in\R$ zbiór
    $\{x\in X:f(x)\geq a\}$ należy do $\mathcal{F}$;
  \item dla każdego $a\in\R$ zbiór
    $\{x\in X:f(x)<a\}$ należy do $\mathcal{F}$;
  \item dla każdego $a\in\R$ zbiór
    $\{x\in X:f(x)\leq a\}=f^{-1}((-\infty,a])$ należy do $\mathcal{F}$.
\end{enumerate}
\end{twierdzenie}

\begin{proof}
\emph{$1\Rightarrow2$.} Dla każdego $a\in\R$ zachodzi
\[
  \{x\in X:f(x)\geq a\}
  =\bigcap_{n=1}^{\infty}\left\{x\in X:f(x)>a-\frac1n\right\}.
\]
Każdy zbiór po prawej należy do $\mathcal{F}$ na mocy warunku 1,
a sigma-ciało jest zamknięte na przeliczalne przekroje.

\emph{$2\Rightarrow3$.} Korzystamy z dopełnienia względem $X$:
\[
  \{x\in X:f(x)<a\}
  =X\setminus\{x\in X:f(x)\geq a\}\in\mathcal{F}.
\]

\emph{$3\Rightarrow4$.} Analogicznie do pierwszego przejścia mamy
\[
  \{x\in X:f(x)\leq a\}
  =\bigcap_{n=1}^{\infty}\left\{x\in X:f(x)<a+\frac1n\right\}
  \in\mathcal{F}.
\]

\emph{$4\Rightarrow1$.} Ponownie bierzemy dopełnienie:
\[
  \{x\in X:f(x)>a\}
  =X\setminus\{x\in X:f(x)\leq a\}\in\mathcal{F}.
\]
\end{proof}

\begin{definicja}[Funkcja mierzalna]
\label{def:funkcja-mierzalna}
Funkcję $f\colon X\to\R$ nazywamy mierzalną względem $\mathcal{F}$,
jeżeli spełnia dowolny, a więc każdy z warunków
twierdzenia~\ref{tw:warunki-mierzalnosci-funkcji}.
\end{definicja}

\begin{uwaga}[Wartości w prostej rozszerzonej]
\label{uw:funkcja-mierzalna-rozszerzona}
Tak samo definiujemy mierzalność funkcji
$f\colon X\to\Rbar=[-\infty,+\infty]$:
sprawdzamy te same zbiory poziomicowe dla wszystkich $a\in\R$.
Dowód równoważności czterech warunków pozostaje taki sam.
\end{uwaga}

\section{Przeciwobrazy i kryterium borelowskie}
\label{sec:przeciwobrazy-borelowskie}

\begin{fakt}[Przedziały domknięte]
\label{fakt:przeciwobraz-domknietego}
Jeżeli $f\colon X\to\R$ jest mierzalna, to dla $a\leq b$
\[
  f^{-1}([a,b])\in\mathcal{F}.
\]
\end{fakt}

\begin{proof}
Przeciwobraz zachowuje przekroje, zatem
\begin{align*}
  f^{-1}([a,b])
    &=f^{-1}\bigl((-\infty,b]\cap[a,+\infty)\bigr)\\
    &=\{x\in X:f(x)\leq b\}\cap\{x\in X:f(x)\geq a\}
      \in\mathcal{F}.
\end{align*}
\end{proof}

\begin{wniosek}[Przeciwobraz punktu]
\label{wn:przeciwobraz-punktu}
Dla funkcji mierzalnej $f$ i każdego $a\in\R$ mamy
$f^{-1}(\{a\})\in\mathcal{F}$.
\end{wniosek}

\begin{proof}
Wystarczy wziąć $a=b$ w fakcie~\ref{fakt:przeciwobraz-domknietego}.
\end{proof}

\begin{uwaga}
Dla mierzalnej funkcji $f\colon X\to\Rbar$ mierzalne są również
przeciwobrazy wartości nieskończonych. W szczególności
\[
  f^{-1}(\{+\infty\})
  =\bigcap_{n=1}^{\infty}\{x\in X:f(x)\geq n\}\in\mathcal{F}.
\]
Dla $-\infty$ używamy analogicznie przekroju zbiorów
$\{x\in X:f(x)\leq-n\}$.
\end{uwaga}

\begin{fakt}[Pozostałe rodzaje przedziałów]
\label{fakt:przeciwobrazy-przedzialow}
Dla funkcji mierzalnej $f\colon X\to\R$ i $a<b$ zbiory
\[
  f^{-1}((a,b)),\qquad f^{-1}((a,b]),\qquad f^{-1}([a,b))
\]
należą do $\mathcal{F}$.
\end{fakt}

\begin{proof}
Każdy z tych przedziałów otrzymujemy z $[a,b]$, usuwając odpowiednie
końce. Przykładowo
\[
  f^{-1}((a,b])=f^{-1}([a,b])\setminus f^{-1}(\{a\})
  \in\mathcal{F}.
\]
Dla pozostałych dwóch przypadków usuwamy odpowiednio oba końce
albo tylko $b$.
\end{proof}

\begin{twierdzenie}[Kryterium borelowskie]
\label{tw:mierzalnosc-borelowska}
Funkcja $f\colon X\to\R$ jest mierzalna wtedy i tylko wtedy,
gdy dla każdego $B\in\mathcal{B}(\R)$ zachodzi
\[
  f^{-1}(B)\in\mathcal{F}.
\]
\end{twierdzenie}

\begin{proof}
\emph{$\Leftarrow$.} Każda półprosta $(a,+\infty)$ jest zbiorem
borelowskim, zatem spełniony jest warunek 1
twierdzenia~\ref{tw:warunki-mierzalnosci-funkcji}.

\emph{$\Rightarrow$.} Załóżmy, że $f$ jest mierzalna, i połóżmy
\[
  \mathcal{A}=\{B\subseteq\R:f^{-1}(B)\in\mathcal{F}\}.
\]
Rodzina $\mathcal{A}$ jest sigma-ciałem. Istotnie,
$f^{-1}(\varnothing)=\varnothing\in\mathcal{F}$, a tożsamości
\[
  f^{-1}(\R\setminus B)=X\setminus f^{-1}(B),\qquad
  f^{-1}\left(\bigcup_{n=1}^{\infty}B_n\right)
  =\bigcup_{n=1}^{\infty}f^{-1}(B_n)
\]
zapewniają zamknięcie na dopełnienia i przeliczalne sumy.
Z mierzalności $f$ wynika, że $(a,+\infty)\in\mathcal{A}$ dla każdego
$a\in\R$. Półproste te generują $\mathcal{B}(\R)$, więc
$\mathcal{B}(\R)\subseteq\mathcal{A}$. Jest to dokładnie żądany warunek.
\end{proof}

\section{Działania na funkcjach mierzalnych}
\label{sec:dzialania-na-mierzalnych}

\begin{twierdzenie}
\label{tw:dzialania-na-mierzalnych}
Jeżeli $f,g\colon X\to\R$ są mierzalne i $c\in\R$, to funkcje
\[
  f+c,\qquad cf,\qquad f+g,\qquad f-g,\qquad fg
\]
są mierzalne. Iloraz $f/g$ jest mierzalny na zbiorze
$D=\{x\in X:g(x)\neq0\}$ z sigma-ciałem
$\mathcal{F}|_D=\{A\cap D:A\in\mathcal{F}\}$.
\end{twierdzenie}

\begin{proof}[Dowód dla przesunięcia i mnożenia przez stałą]
Ustalmy $a\in\R$. Dla przesunięcia mamy
\[
  \{x\in X:f(x)+c>a\}
  =\{x\in X:f(x)>a-c\}\in\mathcal{F}.
\]
Jeżeli $c>0$, to
\[
  \{x\in X:cf(x)>a\}
  =\{x\in X:f(x)>a/c\}\in\mathcal{F}.
\]
Jeżeli $c<0$, to po podzieleniu nierówności zmienia się jej zwrot:
\[
  \{x\in X:cf(x)>a\}
  =\{x\in X:f(x)<a/c\}\in\mathcal{F}.
\]
Wreszcie, dla $c=0$ otrzymujemy
\[
  \{x\in X:cf(x)>a\}
  =\begin{cases}
    \varnothing, & a\geq0,\\
    X, & a<0.
  \end{cases}
\]
Oba zbiory należą do $\mathcal{F}$.
\end{proof}
